\selectlanguage{english}
\begin{abstract}
The photoelectric effect experiment was conducted using the voltmeter-ammeter method, in which the characteristic curve was obtained for two filters, blue and purple, with three different intensities for each filter, aiming to determine the stopping potential for each. Subsequently, the Planck constant was calculated through a linear relationship between the kinetic energy and the frequencies corresponding to each filter, yielding a result of \( h = \qty{4.00}{\eV\s} \), with a relative percentage error of \qty{3}{\percent} compared to the value reported in the literature.

\end{abstract}

\begin{IEEEkeywords}
Photoelectric effect, Planck constant, Work function, Characteristic curves.
\end{IEEEkeywords}

\selectlanguage{spanish}
\begin{abstract}
Se realizó el experimento del efecto fotoeléctrico utilizando el método voltímetro-amperímetro, en el cual se obtuvo la curva característica para dos filtros de color azul y morado, con tres intensidades diferentes en cada filtro, con el objetivo de determinar el potencial de frenado para cada uno. Posteriormente, se calculó la constante de Planck mediante una relación lineal entre la energía cinética y las frecuencias correspondientes a cada filtro, obteniendo como resultado \( h = \qty{4.00}{\eV\s} \), con un error relativo porcentual de \qty{3}{\percent} respecto al valor registrado en la literatura.
\end{abstract}
Efecto fotoeléctrico, constante de Planck, función de trabajo, curvas características.
\begin{IEEEkeywords}

\end{IEEEkeywords}
